\documentclass[11pt]{article}
\usepackage[letterpaper,margin=0.85in]{geometry}
\usepackage{amsmath,amssymb}
\usepackage{booktabs,longtable,tabularx,array}
\usepackage{pdflscape}
\usepackage{xcolor}
\usepackage[colorlinks=true,allcolors=blue]{hyperref}
\usepackage[round,authoryear]{natbib}
\usepackage{enumitem}
\setlist{nosep,leftmargin=*}
\newcolumntype{Y}{>{\raggedright\arraybackslash}X}
\newcommand{\Lean}{\textsc{Lean}}
\newcommand{\Yes}{yes}
\newcommand{\Partial}{partial}
\newcommand{\No}{no}

\title{Semantic Alignment in LLM-Assisted Formalization:\\Detailed Related-Work and Tooling Notes}
\author{Internal reference for the four-page workshop paper}
\date{Search updated 4 September 2026}

\begin{document}
\maketitle

\begin{abstract}
This memo surveys work related to semantic correspondence between informal
mathematics and formal statements, with emphasis on Lean and on workflows that
maintain a formalization over time.  It is intentionally broader than the
related-work section of the workshop submission.  The survey separates
reference-based formal equivalence, learned or LLM-based source-to-formal
judging, compiler-grounded semantic tests, proof-structure preservation,
long-horizon source review, and human review.  The final sections identify the
closest systems for semantic review and organizes them as a practical map for
researchers entering formalization through LLM-assisted tools.  It also records
which claims about our local dashboard are supported by the literature search as
of the date above.
\end{abstract}

\tableofcontents

\section{Scope and search procedure}

The workshop paper studies a practical question: after Lean accepts a theorem,
what evidence supports the additional judgment that the theorem represents a
particular source result?  The search therefore included work on statement
autoformalization, semantic-equivalence metrics, critics, full theorem-and-proof
translation, source-linked project formalization, independent review, blueprint
maintenance, and human-in-the-loop formalization.  Work devoted only to proof
search from an already fixed formal theorem is included only when it contributes
a semantic-evaluation mechanism.

The literature search was refreshed on 4 September 2026.  Queries combined
terms such as \emph{autoformalization}, \emph{semantic alignment},
\emph{semantic fidelity}, \emph{faithfulness}, \emph{equivalence},
\emph{source review}, \emph{Lean}, \emph{blueprint}, \emph{adversarial review},
and \emph{long-horizon formalization}.  Primary records were preferred from
NeurIPS, ICLR/OpenReview, ACL/EMNLP, and arXiv.  The search also followed
citations and related-work sections of the closest papers.  Very recent arXiv
preprints are labeled as such in the discussion.  Terminology in this area is
changing rapidly, so a later search can reveal work that uses a different name
for the same problem.

\section{A taxonomy of semantic-alignment evidence}

The papers fall into several overlapping groups.

\begin{enumerate}
  \item \textbf{Reference-based formal equivalence.} A generated formal
  statement is compared with a trusted formal statement.  Bidirectional
  provability can provide strong evidence when the two statements live in a
  compatible environment.  BEq/BEq+ and ProofBridge are representative.

  \item \textbf{Source-to-formal semantic judging.} A model directly judges an
  informal/formal pair, often with human-labeled positive and negative examples
  or human calibration.  FormalAlign, CriticLean, FormalMATH, and Beyond
  Compilation occupy this space.

  \item \textbf{Compiler-grounded semantic tests without one exact reference
  theorem.}  Theorem Testing uses downstream dependent theorems as integration
  tests.  ShadowBench represents the intended theorem by a complete family of
  auxiliary ``shadow'' statements and checks both implication directions.

  \item \textbf{Roundtrip and consistency methods.}  A formalization is
  translated back, regenerated, or compared with independently generated
  candidates.  These methods provide semantic evidence when a trusted formal
  reference is unavailable, although their guarantees depend on the auxiliary
  translation or comparison mechanism.

  \item \textbf{Proof-structure correspondence.}  ProofFlow and related work
  evaluate preservation of the structure of an informal proof in addition to
  theorem-statement meaning.

  \item \textbf{Long-horizon source review and maintenance.}  LeanMarathon,
  LeanFlow, Beyond the Library, FormaTheoria, M2F, Automatic Textbook
  Formalization, and LeanArchitect address larger artifacts where source
  context, dependencies, interfaces, and later edits affect correspondence.

  \item \textbf{Human and expert review.}  Several projects use domain experts
  to establish or calibrate semantic correspondence.  This includes Beyond the
  Library, FormalScience, the expert-review study of Ilin and Nugent, and the
  human-workflow study of \citet{collins2026humanai}.
\end{enumerate}

These groups answer different questions.  Equivalence to a trusted formal
reference is useful for benchmark scoring.  A source-driven research project can
lack such a reference, can reconstruct assumptions from several passages, and
can change both its source model and Lean representation during development.
For that setting, semantic review becomes a maintained scholarly record as well
as an evaluation decision.

\begin{landscape}
\section{Comparison table}
\small
\begin{longtable}{@{}p{0.18\linewidth}p{0.12\linewidth}p{0.20\linewidth}p{0.18\linewidth}p{0.21\linewidth}@{}}
\caption{Related approaches ordered roughly by the kind of semantic evidence they provide.  ``Reference'' means the object used to define or evaluate the intended meaning; it does not imply that every system requires a gold Lean theorem.}\label{tab:rw-matrix}\\
\toprule
Work & Unit & Semantic mechanism & Reference / source & Relevance to the dashboard \\
\midrule
\endfirsthead
\toprule
Work & Unit & Semantic mechanism & Reference / source & Relevance to the dashboard \\
\midrule
\endhead
\bottomrule
\endfoot

Wu et al. \citep{wu2022autoformalization} & statement & human judgment / exact translation evaluation & informal problem paired with formal specification & Establishes LLM autoformalization as a translation task and an early semantic target. \\

ProofNet \citep{azerbayev2023proofnet} & statement & benchmark pair and manual correctness & NL theorem and Lean theorem & Provides the paired-statement benchmark paradigm used by later alignment metrics. \\

FormalAlign \citep{lu2025formalalign} & statement & learned informal--formal alignment score using generation and contrastive representation losses & paired NL/formal data plus synthetic misalignments & Direct automated source-to-formal evaluator; does not itself maintain a long-lived review record. \\

BEq \citep{liu2025beq} & statement & bidirectional extended definitional equivalence & trusted formal reference & Compiler/prover-grounded equivalence metric; strongest when a reference theorem already exists. \\

BEq+ / ProofNetVerif \citep{poiroux2025reliable} & statement & deterministic bidirectional proof search; human-labeled metric benchmark & trusted formal reference plus NL equivalence labels & Provides a reproducible semantic metric and research-level RLM25 benchmark. \\

Symbolic equivalence + semantic consistency \citep{li2024symbolic} & statement & equivalence among candidates plus informalization/embedding consistency & original NL statement and generated candidate set & Uses consistency to select among several formalizations rather than preserving reviewer evidence over time. \\

FormalMATH \citep{yu2025formalmath} & statement & multi-LLM semantic verification plus disproof filtering and human verification & original problem statement & Demonstrates a layered filtering pipeline before expert acceptance. \\

CriticLean \citep{peng2025criticlean} & statement & trained semantic critic; critic-guided RL & informal statement with labeled formal candidates & Closest learned-critic analogue to the skeptical-review role. \\

Beyond Compilation \citep{zhang2026beyondcompilation} & statement & cross-model semantic judging calibrated with human audits & original NL statement & Quantifies the compile-pass / judged-faithfulness gap and separates validity from faithfulness. \\

Theorem Testing / T$^2$ \citep{kim2026theoremtesting} & theorem in repository & compile downstream successor theorems as semantic tests & existing dependency graph & Uses repository context as automatically extracted semantic constraints; weak for standalone/frontier theorems. \\

ShadowBench / SA-Pass \citep{han2026shadowbench} & statement & candidate implies all shadows and conjunction of shadows implies candidate & expert-reviewed complete shadow set derived from intended theorem & Very close in clause decomposition and bidirectional checking; optimized for automatic benchmark evaluation rather than maintenance of source-review provenance. \\

Roundtrip verification \citep{amrollahi2026roundtrip} & specification & NL$\rightarrow$formal$\rightarrow$NL$\rightarrow$formal, then formal equivalence & original NL statement; second formalization generated through roundtrip & Reference-free semantic evidence and targeted repair; demonstrated on traffic rules rather than Lean mathematics. \\

ProofBridge \citep{jana2026proofbridge} & theorem + proof & joint NL/FL embeddings plus bidirectional Lean equivalence for semantic correctness & paired NL/Lean theorem-proof data & Extends semantic evaluation to whole proofs and uses formal equivalence as a metric. \\

ProofFlow \citep{cabral2026proofflow} & theorem + proof & dependency DAG and ProofScore combining syntax, semantics, structure & manually annotated proof steps and dependency graphs & Adds structural fidelity beyond theorem correspondence. \\

FormalScience \citep{meadows2026formalscience} & scientific statement + proof & human-in-the-loop review and taxonomy of semantic drift & domain-expert interpretation of physics source & Shows domain-object collapse and abstraction changes that can remain formally valid. \\

M2F \citep{wang2026m2f} & document/project & decomposition into atomic blocks; dependency-ordered statement compilation; verifier feedback & long-form source document & Project-scale context and fixed-signature proof stage; adjacent to source-linked maintenance. \\

Automatic Textbook Formalization \citep{gloeckle2026textbook} & textbook/project & large multi-agent development with side-by-side blueprint & graduate textbook & Demonstrates scale and a human-facing informal/formal blueprint. \\

Beyond the Library \citep{moakhar2026beyondlibrary} & research paper & faithfulness judge plus human-expert validation; source-linked public formalizations & target research papers & Closest source-driven research-paper setting with human validation. \\

LeanFlow \citep{milikic2026leanflow} & paper/project & statement/source review and workflow mechanisms; BEq+ calibration & mathematical papers and RLM25 slices & Studies auditability and workflow effects at document scale. \\

LeanMarathon \citep{zhang2026leanmarathon} & research development & adversarial target-review loop before proof DAG execution & shared blueprint pairing formal types with natural-language targets & Close analogue of producer/reviewer separation and target-fidelity stabilization. \\

LeanArchitect \citep{zhu2026leanarchitect} & project blueprint & Lean annotations generate synchronized blueprint metadata and dependency information & Lean declarations linked to informal blueprint nodes & Closest maintenance tool for keeping informal/formal project metadata synchronized. \\

FormaTheoria \citep{nie2026formatheoria} & distributed literature / theory & independent semantic review, provenance, protected declarations, source reconciliation, review gates & section-level context from many mathematical sources & Strongest long-horizon neighbor: explicit provenance and independent review with immutable approved state. \\

Ilin--Nugent \citep{ilin2026sorries} & project/library contribution & expert review after a no-sorry compiling formalization & expert expectations for definitions, generality, organization, API & Evidence that compilation can leave higher-level formalization defects requiring expert review. \\

\end{longtable}
\end{landscape}

\section{Statement-level semantic evaluation}

\subsection{From autoformalization accuracy to semantic equivalence}

\citet{wu2022autoformalization} showed that LLMs could translate a nontrivial
fraction of competition mathematics into formal specifications, reporting
25.3\% perfect Isabelle/HOL translations.  ProofNet then supplied 371
undergraduate examples containing an informal theorem, Lean theorem, and
informal proof \citep{azerbayev2023proofnet}.  These works established the
paired informal/formal benchmark setting in which later semantic metrics are
defined.

A key transition came from replacing lexical or compile-only scores with
semantic checks.  \citet{liu2025beq} introduced BEq, which asks whether two
formal statements can be shown equivalent using a bidirectional extended
notion of definitional equivalence.  Their Con-NF benchmark contains 961
informal/formal pairs from research mathematics, explicitly moving beyond
self-contained competition problems.  \citet{poiroux2025reliable} refined this
line with BEq+, a deterministic reference-based metric, and introduced
ProofNetVerif with 3,752 human-annotated examples for evaluating semantic
metrics.  Their RLM25 benchmark adds 619 pairs sampled from six research-level
formalization projects.  BEq+ achieved 98.0\% precision and 48.3\% recall on the
reported ProofNetVerif binary evaluation; the paper also documents false
positives and false negatives that arise from the strength and framing of the
proof search.

Reference-based equivalence is especially attractive when a trusted Lean
formalization already exists.  The research setting in the workshop paper often
starts before such a reference exists.  The intended result may need to be
reconstructed from a theorem statement, section assumptions, definitions, and
representation conventions.  The dashboard therefore records the source
analysis that creates the comparison target instead of assuming a pre-existing
formal target.

\subsection{Learned and LLM semantic critics}

FormalAlign directly scores informal/formal alignment.  Its model is trained on
both autoformalization and a contrastive representation-alignment objective,
and evaluation uses synthetic misalignment strategies on FormL4 and miniF2F
\citep{lu2025formalalign}.  The approach turns semantic checking into a learned
selection problem.

CriticLean makes the critic a first-class learned component
\citep{peng2025criticlean}.  It introduces CriticLeanGPT, CriticLeanBench, and a
critic-guided reinforcement-learning loop.  This is conceptually close to the
case study's use of a fresh context whose task is to challenge an earlier
correspondence claim.  The workshop paper's retrospective evidence does not
support a controlled comparison of critic prompting strategies or model
families, so CriticLean should be cited as a complementary systematic treatment
of the critic role.

Beyond Compilation provides one of the clearest measurements of the semantic
gap in current Lean statement generation \citep{zhang2026beyondcompilation}.
On a 400-entry graduate-level benchmark, the full tool-augmented agent reported
89.5\% compilation and 60.5\% consensus faithfulness.  Its evaluation combines
Lean compilation, cross-model semantic judging, and targeted human calibration.
This directly supports separating compiler validity from source correspondence
in the workshop paper.

FormalMATH uses several layers during benchmark construction: a specialized
statement formalizer, multi-LLM semantic verification, negation-based disproof
filtering, and expert verification \citep{yu2025formalmath}.  This shows another
way to allocate expensive human review after automated filters have removed
clear failures.

\subsection{Compiler-grounded semantic tests}

Theorem Testing treats existing theorem dependencies as an integration-test
suite \citep{kim2026theoremtesting}.  A generated theorem passes when dependent
successor theorems continue to compile.  The associated T$^2$ benchmark contains
2,206 theorem-generation problems from five Lean repositories and pairs each
target with downstream uses.  The method derives semantic constraints without
manual reference annotations, but it is least informative for theorems with few
or no downstream dependents.  This is relevant to the dashboard's display of
proof dependencies: dependency information can act as evidence about expected
interfaces even when it does not by itself establish source correspondence.

ShadowBench is the closest current benchmark to clause-level semantic
obligations \citep{han2026shadowbench}.  SA-Pass represents an intended formal
statement by auxiliary ``shadow'' theorems.  Every shadow must follow from the
target, and their conjunction must imply the target, so a complete shadow set is
Lean-checked to characterize the target in both directions.  ShadowBench
contains 178 postgraduate- to research-level problems.  Its reported Claude
Code/Numina configuration compiled 61.8\% of problems but passed SA-Pass on
11.2\%; SA-Pass achieved 98.8\% binary agreement with expert judgments across
six agentic configurations.

The overlap with the dashboard is important.  Both approaches decompose a
semantic target into smaller obligations and use Lean-checkable relations where
possible.  ShadowBench begins from an intended formal target and builds a set of
formal shadows suitable for automatic scoring.  The dashboard begins from a
mathematical source and preserves the reviewer-visible path from source clauses
to elaborated Lean declarations, including assumptions inherited from surrounding
text and explicit bridge theorems when the Lean representation differs.  The
systems could be combined: accepted dashboard clauses could motivate shadow
statements, while SA-Pass-like checks could strengthen clauses that admit a
formal equivalence characterization.

\subsection{Consistency and roundtrip evidence}

\citet{li2024symbolic} selects among several autoformalization candidates using
symbolic equivalence among formal outputs and semantic consistency obtained by
informalizing candidates and comparing them with the original statement.  The
method is useful when generation diversity produces a correct candidate that
pass@1 misses.

\citet{amrollahi2026roundtrip} studies a reference-free roundtrip in a formal
specification domain: natural language is formalized, translated back, and
formalized again; the two formal results are then checked for logical
equivalence.  On 150 traffic rules, diagnosis-guided repair raised formal
equivalence from 45--61\% to 83--85\% for the tested models.  The domain differs
from Lean mathematics, but the method clarifies an important option when no
trusted formal reference is available.

A very recent preprint, \citet{mohammad2026faithfulnessgap}, proposes
bidirectional provability fingerprinting and counterfactual probes for
natural-language/formal semantic equivalence.  It is included here because its
problem statement is directly relevant; its claims should be treated as recent
preprint results until further independent evaluation is available.

\section{Faithfulness of proofs, not only theorem statements}

Statement correspondence is sufficient to know what proposition Lean has
proved, but some applications also ask whether the formal proof represents the
argument in the source.  ProofBridge jointly autoformalizes theorem statements
and proofs and evaluates semantic correctness using bidirectional Lean
equivalence in addition to type correctness \citep{jana2026proofbridge}.
ProofFlow goes further by constructing a dependency DAG from the informal proof
and formalizing steps as intermediate lemmas; its ProofScore combines syntactic,
semantic, and structural fidelity \citep{cabral2026proofflow}.

The workshop dashboard currently focuses on theorem/source correspondence and
on the bridge declarations needed to justify representation changes.  It can
display proof dependencies, but it does not claim that the generated proof
preserves the rhetorical or structural organization of the paper's proof.  This
boundary should be explicit in the submission so that ``semantic alignment'' is
understood at the level the artifact actually reviews.

\section{Long-horizon and literature-driven formalization}

The closest work to the workshop contribution increasingly comes from
project-scale formalization rather than isolated-statement benchmarks.

\subsection{Document-to-project systems}

M2F divides a mathematical document into atomic blocks, infers dependencies,
compiles statement skeletons, and then repairs proofs under fixed signatures
\citep{wang2026m2f}.  It demonstrates that statement stabilization and proof
construction can be separate phases at project scale.

Automatic Textbook Formalization reports a 500+ page graduate algebraic
combinatorics formalization with 130K lines and 5,900 Lean declarations, built
by many collaborating agents; the project also exposes a side-by-side blueprint
site \citep{gloeckle2026textbook}.  The scale makes synchronization between
informal and formal artifacts a practical concern even when the paper's primary
focus is automated construction.

LeanFlow explicitly studies workflow choices in document-to-project
formalization \citep{milikic2026leanflow}.  Its evaluation includes project
completion, API calls and token counts, with BEq+ used as complementary
statement-level calibration on RLM25.  This is useful evidence that semantic
metrics and project-scale workflow measurements can coexist without serving the
same purpose.

Beyond the Library targets main theorems and proofs from research papers,
extends missing formal infrastructure, uses a faithfulness judge in the
pipeline, and reports human-expert validation of generated research-paper
formalizations \citep{moakhar2026beyondlibrary}.  It is a close comparison for
the Davis--Kahan case because the source is research literature and the target
requires mathematics outside the existing library.

\subsection{Adversarial target review and protected project state}

LeanMarathon makes target review an explicit first stage.  Its blueprint pairs
formal theorem types with natural-language statements and proof information;
a Target-Reviewer and Refiner run an adversarial loop until the target is
accepted before parallel proof construction begins \citep{zhang2026leanmarathon}.
This is the closest published analogue to the producer/reviewer separation that
emerged in the Davis--Kahan case study.  The workshop paper should therefore
claim retrospective evidence for that pattern in this project, rather than
claiming adversarial target review itself as a new method.

FormaTheoria is an even closer comparison at the literature level
\citep{nie2026formatheoria}.  It coordinates source acquisition, translation,
dependency discovery, proof construction, independent review, provenance, and
cross-source reconciliation.  Approved declarations are protected from
unsupported changes, including kernel-level checks of statement immutability.
The paper reports that eleven of fourteen surveyed sections were returned after
the first semantic review, and its review framework explicitly separates
mechanical checks from independent source review.  FormaTheoria therefore
substantially overlaps with broad claims about long-horizon semantic review,
provenance, review gates, and protecting accepted statements.

The dashboard contribution should be positioned at a narrower layer.  It is a
reviewer-facing evidence artifact for one source-to-Lean correspondence: source
fragments and inherited context, clause-level relations, elaborated Lean
statements, bridge declarations, and hashes that make an earlier acceptance
stale after either side changes.  The current literature search found several
systems that preserve provenance or protect theorem statements, but did not
find another published interface described as binding a human semantic-review
decision to hashes of both the exact source evidence and compiler-elaborated
Lean evidence while presenting the clause mapping in one review view.  This is
a literature-search conclusion, not a priority claim; the submission should
avoid ``first'' language.

\subsection{Blueprint synchronization}

LeanArchitect addresses maintenance from the formal side
\citep{zhu2026leanarchitect}.  Lean annotations associate declarations with
blueprint metadata, dependencies are inferred, and LaTeX blueprint material is
generated from the Lean development.  The system reduces duplication between
informal and formal project representations and reports exposing latent
inconsistencies in existing blueprints.  It is an important comparison for the
dashboard's drift detection.  LeanArchitect synchronizes project metadata from
Lean annotations; the dashboard stores the evidence used to accept a specific
source-to-formal semantic relation and asks for renewed review when pinned
source or statement evidence changes.

\section{Human review and domain semantics}

The expert-review study by \citet{ilin2026sorries} begins from a compiling,
no-sorry formalization of Grothendieck's vanishing theorem.  Expert review found
problems in definitions, theorem generality, file organization, and the API,
after which the authors performed a review-driven refactor.  This study is
broader than theorem/source semantic equivalence, but it shows that formal
acceptance and library-quality mathematical judgment operate at different
levels.

FormalScience studies autoformalization in physics with a human-in-the-loop
pipeline \citep{meadows2026formalscience}.  Its semantic-drift analysis includes
cases where operators become scalars or analytic content is moved into
hypotheses.  Those examples closely resemble the kind of representation and
scope changes that a source-to-Lean review interface needs to expose.

\citet{collins2026humanai} provides complementary human-factors evidence.  Their
mixed-methods study examines what users want from AI formalization tools and how
they actually combine tools during formalization.  The paper is useful context
for the workshop paper's emergent workflow, while the workshop's retrospective
Git analysis has a different goal: reconstructing semantic-review reversals in
one long-running project.


\section{A practical map of current semantic-review tooling}

For a reader beginning an LLM-assisted Lean project, the current landscape can
be organized by the question that needs to be answered.

\begin{itemize}
  \item \textbf{How can I compare an informal theorem with a formal statement?}
  BEq/BEq+, CriticLean, Beyond Compilation, Theorem Testing, and ShadowBench
  provide complementary automated approaches ranging from formal equivalence to
  learned critics and compiler-grounded shadow statements.

  \item \textbf{Which Lean declarations do I need to understand?}  Lean Atlas
  and Lean Compass compute a semantic review cone around selected target
  theorems, pruning proof-only dependencies and exposing the remaining
  declarations in an interactive web viewer \citep{yanahama2026leanatlas}.

  \item \textbf{How can an author review a research-paper translation?}
  EconCSLib is the closest current example of a paper-facing workflow.  Its
  review dashboard places the paper statement and Lean statement together,
  includes an independently generated Lean-to-TeX rendering and LLM judgment,
  and records human match/mismatch decisions; the surrounding workflow also
  includes validation reports and dependency graphs \citep{garg2026econcs}.

  \item \textbf{How can I browse and evaluate a large autoformalized corpus?}
  ATLAS / AutoformBot pairs textbook targets with Lean declarations, produces
  faithfulness, proof-integrity, and code-quality assessments, and provides a
  visualizer for comparing informal/formal statements and inspecting dependency
  graphs \citep{rammal2026scale}.

  \item \textbf{How should semantic review interact with long-horizon project
  management?}  LeanMarathon, FormaTheoria, and LeanArchitect provide examples
  of adversarial target review, provenance/reconciliation, and synchronization
  between informal project descriptions and Lean declarations
  \citep{zhang2026leanmarathon,nie2026formatheoria,zhu2026leanarchitect}.
\end{itemize}

This map is more useful for the workshop paper than a priority argument.  Our
Davis--Kahan tooling accumulated in one project and now overlaps with several
systems built directly around semantic review.  The paper can use that overlap
to give newcomers concrete places to start and to explain which problems each
approach addresses.

\section{Closest systems for the four-page paper}

For the workshop paper, the most useful comparisons are system-level tools that
a newcomer could plausibly use as part of a semantic-review workflow.

\begin{enumerate}
  \item \textbf{EconCSLib} is the closest paper-facing dashboard comparison.  It
  explicitly separates Lean proof checking from human/LLM review of the
  translation boundary and provides paper-oriented review controls, validation
  reports, and dependency graphs \citep{garg2026econcs}.

  \item \textbf{Lean Atlas / Lean Compass} addresses the complementary problem
  of deciding which project declarations can affect the semantics of selected
  target theorems.  Its interactive viewer and confidence metadata make it a
  direct semantic-verification environment \citep{yanahama2026leanatlas}.

  \item \textbf{ATLAS / AutoformBot} provides large-scale source/formal browsing,
  dependency visualization, and automated faithfulness assessment for textbook
  formalization \citep{rammal2026scale}.

  \item \textbf{ShadowBench / SA-Pass} provides the strongest recent automatic
  statement-level comparison in this survey, using auxiliary shadow statements
  and bidirectional Lean checks \citep{han2026shadowbench}.

  \item \textbf{LeanMarathon, FormaTheoria, and LeanArchitect} show how semantic
  review, provenance, reconciliation, and informal/formal synchronization can be
  integrated into long-running projects
  \citep{zhang2026leanmarathon,nie2026formatheoria,zhu2026leanarchitect}.
\end{enumerate}

The local Davis--Kahan dashboard should therefore be described as an emergent
case-study artifact.  Its useful comparison points are the combination of source
fragments and inherited context, clause-level correspondence notes,
compiler-elaborated Lean statements, representation bridges, and review
invalidation after source or statement changes.  The literature search does not
support a priority claim for semantic-review dashboards in general.

\section{Implications for claims in the four-page paper}

\subsection{Claims well supported by this search}

\begin{itemize}
  \item Lean compilation establishes validity of the formal declaration but
  leaves a distinct source correspondence problem.  This distinction is now
  explicit across FormalAlign, BEq/BEq+, CriticLean, Beyond Compilation,
  Theorem Testing, FormalScience, and ShadowBench.

  \item Semantic evaluation has several established automatic approaches:
  reference equivalence, learned critics, LLM judges, roundtrip checks,
  dependency tests, and shadow statements.  The paper should present the
  dashboard as complementary infrastructure rather than as the introduction of
  semantic alignment as a research problem.

  \item Independent/adversarial target review already appears in long-horizon
  systems, especially LeanMarathon and FormaTheoria.  The Davis--Kahan history
  can contribute a retrospective case study and the concrete dashboard design.

  \item The strongest artifact-specific claim concerns persistent review
  evidence: the dashboard combines source passages and inherited source
  context, clause-level correspondence records, compiler-elaborated Lean
  statements, bridge theorems for representation changes, and source/statement
  hashes used to invalidate stale acceptances.
\end{itemize}

\subsection{Claims to avoid without additional evidence}

\begin{itemize}
  \item Do not claim to introduce semantic alignment, semantic criticism, or
  adversarial review for autoformalization.
  \item Do not claim that independent LLM reviewers are generally more accurate
  than proof-producing LLMs; the repository history is observational and the
  workflow changed over time.
  \item Do not interpret the 14 review reversals as an LLM error rate.  They are
  dependent transitions in one evolving formalization and include checker and
  correspondence-record defects.
  \item Avoid a priority claim for source provenance, statement immutability,
  blueprints, or project-scale formalization; recent systems already provide
  close variants of each.
  \item Avoid claiming that the dashboard mechanically certifies source
  equivalence.  It preserves mechanical evidence and a mathematical review
  judgment; some bridge obligations can be formally checked, but source meaning
  remains partly interpretive.
\end{itemize}

\section{Open comparisons and useful follow-up experiments}

The literature suggests several experiments that would strengthen a longer
version of the paper without being required for the four-page submission.
First, dashboard-reviewed statements could be scored with BEq+ or SA-Pass when a
suitable formal reference or shadow set can be constructed, allowing comparison
between human review records and automatic semantic metrics.  Second, a fixed
sample of historically reopened Davis--Kahan results could be presented to
producer and critic contexts under controlled prompts to test whether role
separation changes detection rates.  Third, the source/statement pinning
mechanism could be evaluated by replaying historical edits and measuring which
semantic changes each class of pin detects.  Fourth, the dashboard data model
could be compared directly with LeanArchitect blueprints and FormaTheoria's
protected-declaration/review state to determine which metadata can be exchanged
rather than duplicated.

The immediate workshop submission can remain modest: present a practical map
of semantic-review systems, distill lessons from the Davis--Kahan experience,
show the local dashboard as one implementation that emerged from those lessons,
and use Git history only as retrospective evidence for why the workflow changed.

\bibliographystyle{plainnat}
\bibliography{references}

\end{document}
